\phantomsection
\chapter{\lang{Theoretical Background}{Fundamentação Teórica}}
\phantomsection

As redes Peer-to-peer estão na computação como conceito e aplicações desde os
primórdios da ARPANET. Na verdade, a própria ARPANET poderia ser descrita
como uma rede de comunicação entre nodos -- ou ``pares'' -- descentralizada
\cite{crocker1969}, o que se refletiu na própria criação da internet. A
internet como conhecemos hoje é uma comunicação entre múltiplas máquinas,
múltiplos pares, para envio e recebimento de recursos. Longe de ser uma
novidade recente, o conceito é maduro e fundamenta a comunicação entre
sistemas desde a década de 1960.

Nos anos 90 ocorreu uma transição do que se conhecia anteriormente como
ARPANET para a World Wide Web, trazendo com a mudança milhões de usuários
leigos buscando e consumindo informações de uma maneira nunca antes vista na
história da humanidade. Nesse aspecto, um modelo vigorou, dominou e se
consolidou como o padrão desse novo ambiente: o modelo cliente-servidor
\cite{fieldingtaylor2000,tanenbaumvansteen2007}. Para entendermos os
principais motivos do sucesso desse novo modelo, precisamos repensar algumas
características das redes peer-to-peer que precisariam ser adaptadas para
essa nova realidade.

Redes peer-to-peer trocam mensagens para definir qual a versão ``correta''
-- muitas vezes subjetiva -- de um pacote, dado ou informação. Com o advento
da internet, milhões de novas máquinas começaram a processar, enviar e
receber esses pacotes/informações a todo momento. Isso, na época, foi um
grande problema relacionado à descoberta e busca de recursos entre os
nodos, saber onde o recurso ou arquivo estava para fazer o upload ou
download \cite{ripeanu2001}.

Naquela época, muitas residências tinham altas -- para a época -- taxas de
download, mas baixíssimas taxas de upload. Fazer todas essas máquinas,
algumas vezes rápidas, outras lentas, responderem requisições teria uma
latência e velocidade muito inferior a simplesmente um data center com
links dedicados de alta velocidade, capazes de responder milhares de
clientes passivos \cite{saroiu2002}.

Somado a isso, os computadores pessoais não precisavam de IP público
estático, portas abertas e nem processar continuamente, garantindo a
simplicidade do modelo no lado do cliente \cite{saltzer1984}.

A adoção do modelo cliente-servidor viabilizou a difusão rápida da Web em
seus primeiros anos. Entretanto, no final da década de 1990, uma nova
tendência surgiu: a explosão de mídias digitais, impulsionada pela
popularização dos arquivos de áudio em MP3. A procura e compartilhamento de
músicas se espalhou e trouxe consigo um aumento no número de bytes
trafegados. A rede não estava mais apenas baixando e fazendo upload de HTML
e arquivos de textos de poucos kbytes, os arquivos de músicas elevaram esse
patamar para os megabytes, causando a explosão das contas das
infraestruturas dos servidores centrais \cite{saroiu2002}. Para ilustrar
esse gargalo de largura de banda (egress bandwidth), em um cenário
hipotético no qual um milhão de usuários querem baixar um HTML de 20kbytes
o servidor precisará lidar com cerca de 20 Gigabytes de dados de saída, com
um arquivo MP3 de 5MB no mesmo servidor, seria necessária uma demanda
massiva de 5 Terabytes, um volume desproporcional para a infraestrutura da
época.

O volume de dados massivo forçava os servidores centrais a um limite o
qual eles não eram capazes de suportar sem encarecer demais o serviço. A
centralização do estado e dos recursos fez com que a disponibilidade de
todo o sistema ficasse refém do servidor central, um único ponto de falha
(Single Point of Failure, ou SPOF) \cite{tanenbaumvansteen2007}. Isso gerou
flash crowds (picos de tráfego) que derrubaram servidores inteiros por
exaustão de CPU, indisponibilizando o serviço para todos os usuários
\cite{padmanabhan2002}.

Os pontos negativos do modelo cliente-servidor vieram à tona, escalabilidade,
disponibilidade, custo e manutenção eram apenas alguns deles. A crise criou
o cenário perfeito para o ressurgimento das redes Peer-to-peer, que
reemergiam como uma alternativa viável e escalável para o problema, em 1999
surge o Napster.

Criado em 1999 por Shawn Fanning e Sean Parker, o Napster como solução de
software provou a viabilidade de uma implementação P2P em larga escala
\cite{saroiu2002}. Sua implementação foi realizada para resolver um
problema comum da época: a dificuldade de baixar arquivos mp3, como
músicas, de sites estáticos ou canais de IRC da época.

O Napster resolveu o gargalo financeiro e o de largura de banda do modelo
cliente-servidor na distribuição desses arquivos, na época, pesados. Por
exemplo, para uma empresa distribuir um vídeo de 1 Terabyte de dados, seria
necessário 1 Terabyte de egress bandwidth dos data centers. Com o Napster
essa infraestrutura de rede, hardware e conexão é repassada para os
próprios usuários da rede, os nodos ou peers. Os usuários que querem baixar
um determinado arquivo fazem essa busca pela rede até encontrar alguém o
servindo. A partir disso, cria-se uma conexão entre os dois, onde os
recursos computacionais de ambos são usados para o compartilhamento,
reduzindo, e muito, o estresse das máquinas centrais dos data centers
\cite{adarhuberman2000}.

Mesmo com tudo isso, o Napster ainda não era uma rede totalmente
descentralizada. Ao invés disso, atuava como uma arquitetura híbrida, ou
P2P de primeira geração, combinando elementos de centralização com
transferência direta de dados \cite{schollmeier2001}. Na execução do
programa, o cliente se conecta com servidores centrais mantidos pelo
Napster. O cliente, então, envia a lista de todos os MP3 armazenados para
compartilhamento, endereço IP e porta de conexão. Os servidores centrais do
Napster mantinham em memória índices globais e dinâmicos em tempo real,
mapeando milhões de arquivos com respectivos endereços de IP. Quando um
usuário procurava por uma música, pelo nome, por exemplo, essa pesquisa era
convertida em uma query e enviada para o servidor central. O servidor
pesquisava em suas tabelas pelo índice e retornava uma lista contando IPs
de usuários que estavam provendo aquele conteúdo e suas velocidades de
conexão (upload e download). Assim que o usuário escolhia a fonte, a
conexão entre ele e o servidor central acabava e começava a conexão TCP
direta entre os nodos/peers. O arquivo requerido pelo usuário consumidor
era, então, enviado do usuário provedor para o seu computador. Embora o
problema do gargalo tenha sido resolvido, dois pontos fracos principais
finalizaram as operações do Napster em 2001:

\begin{itemize}
    \item Caso a infraestrutura dos servidores centrais, detentores das
    tabelas de índices, fosse derrubada, toda a rede como um todo parava de
    funcionar. Algo análogo a um motorista em uma cidade desconhecida com
    GPS desligado. Os servidores centrais eram o novo ponto de falha única
    (SPOF -- Single point of failure) do sistema \cite{tanenbaumvansteen2007};

    \item Esse ponto único de falha foi justamente o foco principal do
    ataque contra o Napster realizado por cortes e juízes americanos, que
    definiram a derrubada desses servidores. Com os servidores fora do ar,
    a rede perdia o seu propósito.
\end{itemize}

Após a queda do Napster, uma rede P2P de primeira geração (híbridas com
servidor central), como ficou conhecida posteriormente, surgiram protocolos
100\% descentralizados e não estruturados, como o Gnutella
\cite{ripeanu2001}. Redes P2P de segunda geração são marcadas,
principalmente, pela eliminação do servidor central usando flooding de
mensagens, não existia servidor central nem tabela de índice. Quando um
peer quer um arquivo, ele envia mensagens diretamente para seus nodos
vizinhos, que fazem o mesmo, consecutivamente, até o nodo detentor do
arquivo ser encontrado. Uma vez encontrado, uma conexão direta é
estabelecida entre ambos para upload ou download do arquivo, a depender da
ponta. O claro gargalo é o flooding, ou inundação, de mensagens na rede,
gerando um tráfego exponencial $O(d^k)$, onde $d$ é o número de conexões e
$k$ o TTL. Em uma rede com milhões de nós a internet dos usuários poderia
ser consumida apenas com tráfego de terceiros.

Almejando o controle do flooding, foram introduzidos o conceito de super
nós, uma arquitetura semi-híbrida \cite{saroiu2002}. Esses super nós faziam
a busca (o flooding) entre eles, que possuíam muitos IPs dinâmicos usados
para direcionar o fluxo otimizado da mensagem. Como exemplo, pode-se
imaginar uma carta enviada pelos correios. Ao invés de cada vizinho
perguntar para os do lado se aquela é a casa correta, apenas as primeiras
casas de cada rua se comunicam, e elas sabem quais casas têm em suas ruas.
No momento em que uma dessas primeiras casas avisa que conhece a casa alvo,
a primeira casa conta diretamente o endereço dela para o envio da carta,
criando-se um caminho muito mais direto até achar a casa correta.

As redes P2P de terceira geração, criadas por Bram Cohen em 2001, foram
popularizadas por clientes de soluções como utorrent \cite{cohen2003}. O
bittorrent revolucionou a transferência de dados ao abandonar o modelo
vigente, o qual enviava arquivos inteiros de um nodo a outro. A
fragmentação em blocos, ou pieces, é definida pela divisão em pequenos
blocos, um usuário não precisa esperar o download terminar para começar a
enviar dados. Enquanto baixa o bloco B de usuário 1, envia o bloco A que já
detém para o usuário 2.

A verdadeira descentralização do bittorrent veio à tona somente com o
surgimento das DHTs, ou distributed hash tables, em especial o protocolo
Kademlia \cite{maymounkovmazieres2002}, que permitiu a operação sem
rastreamento do bittorrent. Em vez de depender de super nós ou inundações
cegas, o Kademlia organiza a rede atribuindo um identificador numérico
único (hash de 160 bits) tanto para os nós quanto para os arquivos. A
localização de um dado é feita calculando uma ``distância lógica''
matemática entre esses identificadores por meio da operação XOR. Para
visualizar o funcionamento de uma DHT, pode-se imaginar um grande catálogo
telefônico global dividido em seções. Um usuário não precisa guardar o
contato de milhões de pessoas na internet; ele guarda apenas uma lista
telefônica reduzida, contendo contatos próximos do seu próprio número e
alguns contatos estratégicos para regiões mais distantes da rede. Quando
deseja buscar um arquivo específico, o nó envia a requisição para o
contato de sua tabela que possui a numeração mais próxima da chave do
arquivo desejado. Esse contato, por sua vez, redireciona o pedido para
alguém ainda mais próximo, afunilando a busca rapidamente até alcançar o
detentor do dado. Esse roteamento estruturado garante que a busca por um
arquivo em uma rede com milhões de nós seja concluída em apenas $O(\log N)$
saltos de mensagens, eliminando completamente a tempestade de tráfego
(flooding) característica das redes de segunda geração.

As distributed hash tables na terceira geração das redes P2P foram um
incontestável sucesso. A distribuição massiva de dados descentralizados em
blocos entre múltiplos milhões de usuários foi alcançada depois de
inúmeros avanços e superações de problemas relacionados a hardware, SO e
redes de computadores \cite{maymounkovmazieres2002,cohen2003}. Entretanto,
essas grandes redes foram projetadas sob a premissa de uma imutabilidade de
dados, um arquivo MP3, ou um vídeo, não mudam de tempos em tempos, seu hash
sempre será aquele, a depender do algoritmo utilizado. Um outro hash
representa um arquivo completamente diferente, sempre. Por outro lado, os
sistemas interativos modernos não são imutáveis, eles mudam rapidamente e
repetidamente, e exigem a alteração constante e simultânea do estado
daquele arquivo, board, ferramenta por múltiplos usuários
\cite{kleppmann2019}.

Em uma arquitetura P2P descentralizada, órfã de um servidor geral agindo
como uma fonte única da verdade (Single Source of Truth), a sincronização
de dados em tempo real enfrenta o dilema da divergência de dados. Uma dada
informação, alterada por dois usuários ao mesmo tempo, precisa ser
sincronizada não só por ambos, mas também por toda a rede, que também pode
estar alterando outros dados paralelamente. A latência e a ordenação dos
fatos podem levar nodos distintos a terem versões diferentes da realidade,
o que pode ser contornado com a reintrodução de um servidor central
\cite{lamport1978,tanenbaumvansteen2007}. O problema é o custo com isso
atrelado, isto é a reintrodução, também, de um ponto único de falha
(Single Point Of Failure).

Como uma alternativa para superar essas limitações sem recorrer ao modelo
cliente-servidor, este trabalho propõe o desenvolvimento e análise de um
Motor de Sincronização Descentralizada para Aplicações em Tempo Real. Para
fundamentar e resolver o conflito de divergência sem autoridade central, o
motor utiliza estruturas de dados específicas para consistência eventual,
como Conflict-free Replicated Data Types (CRDTs) \cite{shapiro2011}. Para
validar e avaliar o desempenho da solução em termos de latência,
resiliência e concorrência, é desenvolvida uma Proof of Concept (PoC)
focada na sincronização espacial de mapas interativos entre múltiplos
usuários simultâneos.